\def\tituloconcepto{Entrelazamiento: el estado de Bell}
% Si el capitulo se compila solo (sin main.tex), esta cabecera abre el documento.
% Antes de % Si el capitulo se compila solo (sin main.tex), esta cabecera abre el documento.
% Antes de % Si el capitulo se compila solo (sin main.tex), esta cabecera abre el documento.
% Antes de \input{cabecera} cada capitulo define \tituloconcepto.
\ifdefined\enmain\else
  \documentclass[11pt,a4paper]{article}
  \input{preambulo}
  \begin{document}
  \title{\tituloconcepto}
  \author{\autor}
  \date{\today}
  \maketitle
\fi
 cada capitulo define \tituloconcepto.
\ifdefined\enmain\else
  \documentclass[11pt,a4paper]{article}
  % Preambulo compartido por todos los documentos. Compilar con XeLaTeX/LuaLaTeX
% (o Tectonic). Si se usa Overleaf, elegir el compilador XeLaTeX.
\usepackage[spanish,es-noshorthands]{babel}
\usepackage{amsmath,amssymb}
\usepackage{booktabs,array}
\usepackage[a4paper,margin=2.5cm]{geometry}
\usepackage{xcolor}
\usepackage{listings}
\usepackage[most]{tcolorbox}
\usepackage{quantikz}
\usepackage[hidelinks]{hyperref}

% ---- Autor (editar a gusto) -------------------------------------------------
\newcommand{\autor}{Genesis Silesky}

% ---- Notacion de Dirac y atajos ---------------------------------------------
% quantikz ya define \ket y \bra; DeclareRobustCommand los redefine sin error.
\DeclareRobustCommand{\ket}[1]{\left|#1\right\rangle}
\DeclareRobustCommand{\bra}[1]{\left\langle#1\right|}
\DeclareRobustCommand{\braket}[2]{\left\langle#1\middle|#2\right\rangle}
\newcommand{\C}{\mathbb{C}}
\newcommand{\R}{\mathbb{R}}
\DeclareMathOperator{\Tr}{Tr}

% ---- Cajas de apoyo --------------------------------------------------------
% "concepto": definiciones previas que hacen falta para entender el capitulo.
\newtcolorbox{concepto}[1]{breakable,colback=blue!4,colframe=blue!55!black,
  fonttitle=\bfseries,title={Concepto previo: #1}}
% "nota": aclaraciones y advertencias.
\newtcolorbox{nota}[1]{breakable,colback=orange!6,colframe=orange!70!black,
  fonttitle=\bfseries,title={#1}}

% ---- Estilos de listados ----------------------------------------------------
\lstdefinestyle{py}{language=Python,basicstyle=\ttfamily\footnotesize,
  breaklines=true,postbreak=\mbox{\textcolor{gray}{$\hookrightarrow$}\space},
  numbers=left,numberstyle=\tiny\color{gray},frame=single,
  commentstyle=\color{green!40!black}\itshape,
  keywordstyle=\color{blue!70!black}\bfseries,
  stringstyle=\color{red!60!black},showstringspaces=false,tabsize=4,
  columns=fullflexible,keepspaces=true}
\lstdefinestyle{salida}{basicstyle=\ttfamily\footnotesize,breaklines=true,
  frame=leftline,backgroundcolor=\color{gray!8},columns=fullflexible,
  keepspaces=true}

  \begin{document}
  \title{\tituloconcepto}
  \author{\autor}
  \date{\today}
  \maketitle
\fi
 cada capitulo define \tituloconcepto.
\ifdefined\enmain\else
  \documentclass[11pt,a4paper]{article}
  % Preambulo compartido por todos los documentos. Compilar con XeLaTeX/LuaLaTeX
% (o Tectonic). Si se usa Overleaf, elegir el compilador XeLaTeX.
\usepackage[spanish,es-noshorthands]{babel}
\usepackage{amsmath,amssymb}
\usepackage{booktabs,array}
\usepackage[a4paper,margin=2.5cm]{geometry}
\usepackage{xcolor}
\usepackage{listings}
\usepackage[most]{tcolorbox}
\usepackage{quantikz}
\usepackage[hidelinks]{hyperref}

% ---- Autor (editar a gusto) -------------------------------------------------
\newcommand{\autor}{Genesis Silesky}

% ---- Notacion de Dirac y atajos ---------------------------------------------
% quantikz ya define \ket y \bra; DeclareRobustCommand los redefine sin error.
\DeclareRobustCommand{\ket}[1]{\left|#1\right\rangle}
\DeclareRobustCommand{\bra}[1]{\left\langle#1\right|}
\DeclareRobustCommand{\braket}[2]{\left\langle#1\middle|#2\right\rangle}
\newcommand{\C}{\mathbb{C}}
\newcommand{\R}{\mathbb{R}}
\DeclareMathOperator{\Tr}{Tr}

% ---- Cajas de apoyo --------------------------------------------------------
% "concepto": definiciones previas que hacen falta para entender el capitulo.
\newtcolorbox{concepto}[1]{breakable,colback=blue!4,colframe=blue!55!black,
  fonttitle=\bfseries,title={Concepto previo: #1}}
% "nota": aclaraciones y advertencias.
\newtcolorbox{nota}[1]{breakable,colback=orange!6,colframe=orange!70!black,
  fonttitle=\bfseries,title={#1}}

% ---- Estilos de listados ----------------------------------------------------
\lstdefinestyle{py}{language=Python,basicstyle=\ttfamily\footnotesize,
  breaklines=true,postbreak=\mbox{\textcolor{gray}{$\hookrightarrow$}\space},
  numbers=left,numberstyle=\tiny\color{gray},frame=single,
  commentstyle=\color{green!40!black}\itshape,
  keywordstyle=\color{blue!70!black}\bfseries,
  stringstyle=\color{red!60!black},showstringspaces=false,tabsize=4,
  columns=fullflexible,keepspaces=true}
\lstdefinestyle{salida}{basicstyle=\ttfamily\footnotesize,breaklines=true,
  frame=leftline,backgroundcolor=\color{gray!8},columns=fullflexible,
  keepspaces=true}

  \begin{document}
  \title{\tituloconcepto}
  \author{\autor}
  \date{\today}
  \maketitle
\fi


\section{Entrelazamiento: el estado de Bell}

\subsection{Objetivo}
Construir el estado de Bell $\tfrac1{\sqrt2}(\ket{00}+\ket{11})$, demostrar que
no se factoriza en dos qubits independientes y observar las correlaciones al
medir.

\subsection{Conceptos previos}
\begin{concepto}{producto tensorial, estado reducido y entropía}
Ver sección~\ref{sec:prerreq}. Tres ideas nuevas:
\begin{itemize}
  \item \textbf{Estado producto:} $\ket a\otimes\ket b$; cada qubit tiene su
    propio estado.
  \item \textbf{Matriz densidad $\rho$:} describe estados puros y mezclados;
    para un estado puro $\rho=\ket\psi\bra\psi$.
  \item \textbf{Traza parcial:} ``olvidar'' un subsistema. El estado del qubit
    restante es $\rho_0=\Tr_1\rho$. Si $\rho_0$ es mezclado, el sistema total
    estaba entrelazado.
\end{itemize}
La entropía de von Neumann es $S(\rho)=-\Tr(\rho\log_2\rho)$: vale $0$ para un
estado puro y $1$ bit para un qubit totalmente mezclado ($\rho=I/2$).
\end{concepto}

\subsection{Fundamento teórico}
La puerta CNOT (control $q_0$, objetivo $q_1$) invierte $q_1$ si $q_0=\ket1$.
Escribiendo los estados como $\ket{q_1q_0}$ (orden de Qiskit):
\begin{equation}
  \text{CNOT}\ket{00}=\ket{00},\qquad \text{CNOT}\ket{01}=\ket{11}.
\end{equation}
El circuito es
\begin{center}
\begin{quantikz}
\lstick{$q_0:\ \ket0$} & \gate{H} & \ctrl{1} & \meter{}\\
\lstick{$q_1:\ \ket0$} &          & \targ{}  & \meter{}
\end{quantikz}
\end{center}
y produce $\ket{\Phi^+}=\tfrac1{\sqrt2}(\ket{00}+\ket{11})$.

\paragraph{No es un producto.} Si $\ket{\Phi^+}=(a\ket0+b\ket1)\otimes
(c\ket0+d\ket1)=ac\ket{00}+ad\ket{01}+bc\ket{10}+bd\ket{11}$, habría que tener
$ad=0$ y $bc=0$ pero $ac=bd=\tfrac1{\sqrt2}\neq0$. Si $a=0$ o $d=0$ se
contradice $ac\ne0$ o $bd\ne0$ respectivamente; así que es imposible.

\paragraph{Cuantificarlo.} Para $\ket{\Phi^+}$, $\rho_0=\Tr_1\rho=I/2$, de
modo que $S(\rho_0)=1$ bit: el conjunto está en un estado puro bien definido,
pero cada parte por separado es máximamente incierta.

\paragraph{Correlación sin comunicación.} Al medir, solo salen $00$ o $11$ (cada
uno con probabilidad $\tfrac12$). Cada resultado individual es aleatorio, por
lo que no se puede usar el entrelazamiento para enviar un mensaje
(\emph{no-signaling}): no hay transmisión de información más rápida que la
luz.

\begin{nota}{Orden de bits}
Qiskit escribe $\ket{q_1q_0}$. Para este estado el orden no importa porque es
simétrico.
\end{nota}

\subsection{Código (completamente comentado)}
\lstinputlisting[style=py]{../05b_entrelazamiento_bell.py}

\subsection{Explicación celda por celda}
\begin{description}
  \item[Celda 2] Construye el circuito y obtiene el vector de estado:
    $(0.707,0,0,0.707)$ para $\ket{00},\ket{01},\ket{10},\ket{11}$.
  \item[Celda 3] \texttt{partial\_trace(sv, [1])} descarta el qubit $1$ y
    \texttt{entropy} calcula $S$: sale $1.0$.
  \item[Celda 4] Copia el circuito, añade mediciones con
    \texttt{measure\_all()} y ejecuta 1000 shots.
\end{description}

\subsection{Salida obtenida}
\lstinputlisting[style=salida]{../salidas/05b_entrelazamiento_bell.txt}
Los conteos exactos de $00$ y $11$ varían entre ejecuciones, pero ambos rondan
$500$ y nunca aparecen $01$ ni $10$.

\subsection{Interpretación}
Los resultados de ambos qubits están perfectamente correlacionados aunque cada
uno por separado sea aleatorio. Es un recurso que no tiene análogo en
probabilidad clásica con variables locales ocultas (desigualdades de Bell).

\subsection{Ejercicios}
\begin{enumerate}
  \item Construya los otros tres estados de Bell
    ($\ket{\Phi^-},\ket{\Psi^\pm}$) agregando puertas $X$ y $Z$ al inicio.
  \item Calcule la entropía de un estado producto (p.\,ej.\ solo $H$ en $q_0$).
  \item Investigue la desigualdad CHSH e impleméntela.
\end{enumerate}

\ifdefined\enmain\else
  \end{document}
\fi

